\documentclass[11pt]{article}
\usepackage{series}
\usepackage{booktabs}
\usepackage{array}
\usepackage{tabularx}

\newcolumntype{L}[1]{>{\raggedright\arraybackslash}p{#1}}
\renewenvironment{abstract}
  {\begin{center}\textbf{Abstract}\end{center}\begin{quote}\small}
  {\end{quote}}

\title{\textbf{Why General Relativity Appears in Perturbative String Theory:}\\
\large A careful MTT interpretation of worldsheet consistency and target-space dynamics}
\author{Peter Nero}
\date{July 2026, Version 2}

\begin{document}
\maketitle

\begin{abstract}
The statement that General Relativity ``falls out of string theory'' is
memorable but easy to misunderstand.  A perturbative string background is
described by a two-dimensional quantum field theory.  Its couplings include
the target-space metric, antisymmetric tensor, dilaton, and, in appropriate
models, gauge and fermionic data.  Requiring quantum Weyl consistency
constrains these couplings.  At leading order those constraints agree, after
the required conventions and field redefinitions, with equations derived
from a target-space effective action containing the Einstein-Hilbert term.
Pure Einstein gravity appears only in a further low-energy and restricted
field corner.  This paper explains that chain without treating measurement,
observation, or mathematical representation as physically privileged acts.
It then gives the restrained Modal Triplet Theory (MTT) interpretation: the
worldsheet and spacetime equations may be two diagnostics emitted by one
preprojection source, but this is established only when explicit source maps
and a controlled comparison square are constructed.  The formal theorem is
owned by the companion technical paper and is not repeated here.  Current
q79 work provides five of twelve worldsheet contract rows, partially provides
two, and leaves five open.  The result is therefore an explanatory bridge and
a concrete completion program, not a proof that GR and string theory are the
same theory or that MTT has already derived either one in full.
\end{abstract}

\section*{Version 2 Revision Note}
\begin{description}[leftmargin=2.8cm,style=nextline]
\item[Supersedes] Version 1.0, released in January 2026.
\item[Reason] The previous version repeated an equivalence theorem, assumed
the relation between MTT proper-time flow and worldsheet RG, and presented
spectral-gap suppression as if it controlled every string and spacetime
correction.
\item[Resolution] Version 2 is an explanatory companion.  It gives the
standard derivation in ordinary language, distinguishes each approximation,
and refers the formal diagnostic-square result to its technical owner.
\item[Retained content] The motivating intuition remains: agreement between
worldsheet and spacetime consistency conditions may point to a shared
upstream description.
\item[Open boundary] The physical q79 worldsheet, its visible--hidden bundle,
the exact infrared conformal theory, GSO and modular data, and the selected
upper MTT action are not yet complete.
\end{description}

\tableofcontents

\section{The phrase and the possible misunderstanding}

Physicists often say that General Relativity falls out of string theory.
What they mean is not that the Einstein equations are hidden in a string and
can be read off without assumptions.  They mean that consistency of the
quantized string worldsheet constrains the background through which the
string propagates, and that the leading target-space equation contains
Einstein's curvature dynamics.

This is a striking result because the calculation begins in two dimensions.
The object being quantized is a map
\[
X:\Sigma\longrightarrow Y
\]
from a worldsheet \(\Sigma\) into a target space \(Y\).  The metric of \(Y\)
appears as a coupling in the two-dimensional theory.  Quantum consistency of
that theory then becomes an equation for the target metric.

The slogan becomes misleading when it erases the intermediate steps.  The
worldsheet is not four-dimensional spacetime.  Its renormalization scale is
not automatically physical time.  The target equations contain more than
the metric.  Their relation to an effective action is perturbative and
scheme-aware.  Four-dimensional GR appears only after a compactification and
low-energy reduction.

\subsection{Four statements of increasing strength}

\begin{center}
\small
\begin{tabularx}{\textwidth}{@{}L{0.27\textwidth}X@{}}
\toprule
Statement & Status\\
\midrule
The metric beta function contains Ricci curvature & Standard leading-order
sigma-model result.\\
Vanishing Weyl-anomaly coefficients agrees with target effective equations &
Standard perturbative result at a declared order and scheme.\\
Pure four-dimensional Einstein gravity follows & Conditional on field
restrictions, compactification, scale separation, and truncation.\\
Worldsheet and spacetime equations descend from one selected MTT source &
An open source-construction problem, formalized conditionally in the
companion technical paper.\\
\bottomrule
\end{tabularx}
\end{center}

Only the first two are part of the ordinary perturbative string derivation.

\section{The worldsheet as an ordinary physical system}

\subsection{Background fields become couplings}

At lowest order, the sigma-model action contains terms of the form
\begin{align}
S_\Sigma[X]
\sim{}&\frac{1}{4\pi\alpha'}\int_\Sigma
\sqrt h\,h^{ab}g_{\mu\nu}(X)
\partial_aX^\mu\partial_bX^\nu\,d^2\sigma \nonumber\\
&+\frac{i}{4\pi\alpha'}\int_\Sigma
\epsilon^{ab}B_{\mu\nu}(X)
\partial_aX^\mu\partial_bX^\nu\,d^2\sigma
+\frac{1}{4\pi}\int_\Sigma\sqrt h\,R^{(2)}\Phi(X)\,d^2\sigma .
\label{eq:sigma}
\end{align}
The target metric \(g\), two-form \(B\), and dilaton \(\Phi\) are therefore
couplings from the worldsheet point of view.  In heterotic theory there are
also gauge-bundle couplings and chiral worldsheet fields.  Ghosts, spin
structures, a path-integral measure, and global data are essential parts of
the quantum theory.

Nothing about this description gives measurement a special ontological role.
The worldsheet theory is a physical quantum system.  A detector is another
physical system coupled to it.  The consistency condition discussed here is
about quantum symmetries and renormalization, not about an observer causing an
outcome.

\subsection{Why Weyl symmetry matters}

The worldsheet metric \(h\) is partly redundant: locally rescaling it should
not change physical string propagation.  Quantization can spoil this Weyl
symmetry.  The failure is encoded in anomaly coefficients, commonly described
through beta functions for the couplings.

For one conventional choice of fields and scheme, the metric coefficient
begins as
\[
\bar\beta^g_{\mu\nu}
=\alpha'\left(
R_{\mu\nu}+2\nabla_\mu\nabla_\nu\Phi
-\frac14H_{\mu\rho\sigma}H_\nu{}^{\rho\sigma}
\right)+O(\alpha'^2).
\]
This formula immediately corrects a common oversimplification.  The equation
is not normally \(R_{\mu\nu}=0\).  It is coupled to the dilaton and the flux
\(H\), and the other background fields have their own equations.

\subsection{A fixed point is not the whole string theory}

Vanishing anomaly coefficients identifies a consistent background in the
perturbative framework.  A complete string construction also asks whether
the worldsheet has the correct central charge, modular and GSO properties,
factorization, anomaly cancellation, infrared limit, and higher-genus
behavior.  A local beta-function calculation can be right while the proposed
global string background is incomplete.

\section{How the target-space action enters}

At string tree level and leading derivative order, the NS--NS sector is
organized by an effective action of the schematic form
\[
S_{\mathrm{eff}}
\sim\int_Y\sqrt{-g}\,e^{-2\Phi}
\left[
R+4|\nabla\Phi|^2-\frac1{12}|H|^2+O(\alpha')
\right]d^Dx .
\]
Its Euler--Lagrange equations match the worldsheet Weyl conditions after
compatible field redefinitions at the calculated order
\cite{CallanEtAl1985,MetsaevTseytlin1987,HullTownsend1988}.

This agreement is the real content behind the slogan.  It is stronger than a
visual resemblance between formulas: two calculations in one perturbative
string construction lead to compatible conditions on the background.  It is
weaker than identity of the worldsheet and target theories.

\subsection{Why field redefinitions matter}

Coordinates on the space of couplings are not unique.  A local
redefinition of \(g\), \(B\), and \(\Phi\) changes the component form of beta
functions and of the effective action.  Comparisons are therefore made in a
fixed scheme or through an explicit redefinition.  A coefficient mismatch
may be a convention mismatch, but that possibility must be demonstrated
rather than assumed.

\subsection{The expansions are separate}

Two perturbative controls are often compressed into one phrase:
\begin{itemize}
\item the alpha-prime expansion controls higher derivatives relative to the
string length; and
\item the string-coupling expansion controls worldsheet genus and quantum
loops in target spacetime.
\end{itemize}
A background can be weakly curved in string units while the string coupling
is not small, or conversely.  Compactification and Kaluza--Klein truncation
introduce further scales.  An MTT projection estimate would be another,
logically separate error.

\section{Where Einstein gravity appears}

The target action contains the Einstein--Hilbert term \(R\), but it also
contains the dilaton, flux, gauge fields, fermions, sources, higher-curvature
terms, and loop corrections.  To reach ordinary four-dimensional GR one
typically needs:
\begin{enumerate}
\item a background whose extra dimensions can be consistently reduced;
\item energies well below the omitted Kaluza--Klein and string modes;
\item a controlled treatment or stabilization of moduli;
\item a specified four-dimensional frame and Newton normalization;
\item a suitable restriction of fluxes, sources, and additional fields; and
\item small or explicitly retained higher-derivative and loop corrections.
\end{enumerate}

In that corner, the leading metric equation takes the familiar form
\[
G_{\mu\nu}+\Lambda g_{\mu\nu}
=8\pi G_{\mathrm N}T_{\mu\nu}+\text{controlled corrections}.
\]
The result is an effective limit.  Calling it effective is not dismissive.
Effective theories can be precise, predictive, and fundamental to the scale
at which their variables are appropriate.

\subsection{Does string theory derive GR?}

There are two reasonable answers, depending on what ``derive'' means.

If it means that perturbative string consistency predicts target equations
whose low-energy metric sector includes Einstein dynamics, then yes.  This is
a standard and important derivation.

If it means that string theory uniquely selects our four-dimensional
background, matter content, couplings, state, and all gravitational
observables, then no such general statement follows from the beta-function
calculation alone.

\section{The restrained MTT interpretation}

MTT asks whether the worldsheet and spacetime descriptions can be understood
as projections of one preprojection object.  This suggestion is attractive
because it turns a relation between lower theories into a question about
commuting maps:
\[
\begin{array}{ccc}
&\text{selected upper source}&\\
\swarrow&&\searrow\\
\text{worldsheet record}&&\text{spacetime record}\\
\downarrow&&\downarrow\\
\text{Weyl diagnostic}&\longleftrightarrow&
\text{effective field equation}.
\end{array}
\]

The diagram is not a proof.  Every arrow needs a definition.  The two lower
records must carry the same source provenance rather than being chosen
independently to agree.  The bottom comparison needs an error estimate and
must not discard an uncontrolled equation.

\subsection{The technical result lives elsewhere}

The companion paper
\emph{Worldsheet and Spacetime Consistency as a Conditional Diagnostic
Square} gives the formal statement.  In brief, if the two diagnostics
intertwine through an explicit comparison map with a bounded inverse on the
compared subspace, then an exact or small worldsheet residual transfers to an
exact or controlled target residual.  It also proves that merely sharing a
fixed point is insufficient.

This paper does not repeat that theorem.  Its job is to explain what the maps
mean physically and why their hypotheses are necessary.

\subsection{What a spectral gap can and cannot do}

A spectral gap may help control leakage from a retained coherent sector.  It
cannot by itself create:
\begin{itemize}
\item the target manifold or its Lorentzian structure;
\item the worldsheet fields, measure, BRST complex, or GSO projection;
\item the sigma-model beta functions;
\item the target effective action and its normalization; or
\item the comparison between worldsheet and target diagnostics.
\end{itemize}
For the same reason, alpha-prime is not automatically the inverse of an MTT
gap.  A source theorem may relate the scales, but the relation must be
derived with dimensions and coefficients intact.

\section{The q79 position}

The current q79 Fu--Yau route is the strongest physical compactification
candidate in the MTT program.  It is not the auxiliary literal product
\(S^1\times\mathrm{Lens}\times\mathrm{Nil}\).  Circle--Lens--Nil language is
best retained as a local rank or operator filtration, while the physical
worldsheet question is posed on the selected q79 geometric branch.

\begin{center}
\small
\begin{tabularx}{\textwidth}{@{}L{0.23\textwidth}X L{0.18\textwidth}@{}}
\toprule
Group & Required content & Status\\
\midrule
Available rows & Time-oriented target, Fu--Yau charge/Bianchi data,
curvature-level visible cancellation, universal heterotic central charge,
and declared low-energy parity limit & 5 of 12\\
Partial rows & Exact infrared conformal theory and modular/GSO
factorization & 2 of 12\\
Open rows & Global differential gerbe, all-orders target, string-field
vertices, infrared/soft completion, and all-genus or nonperturbative
definition & 5 of 12\\
\bottomrule
\end{tabularx}
\end{center}

The authoritative blocker \(B.\mathrm{QG}.01\) requires all twelve rows on
the same physical nonpullback visible--hidden bundle.  The upper-action
blocker \(B.\mathrm{ACTION}.01\) requires the common source whose descendants
produce the accepted lower structures.  Neither is closed by explaining the
standard beta-function result again.

\section{What would count as progress}

The most useful next result would not be another general shadow-bridge
statement.  It would fill a missing object:
\begin{enumerate}
\item construct the physical visible and hidden holomorphic bundles in one
HYM chamber;
\item close the global differential Green--Schwarz data;
\item construct the q79 heterotic worldsheet theory and its infrared SCFT;
\item calculate analytic characters and the modular-invariant GSO sum;
\item derive worldsheet anomaly coefficients and target equations from the
same source at the same order;
\item publish the comparison map, residual, inverse estimate, conventions,
and hashes; and
\item only then compare states and observables.
\end{enumerate}

This sequence is intentionally concrete.  It makes failure informative.
Perhaps the bundle does not exist in the required chamber, the GSO sum fails,
or the comparison becomes singular.  Such an outcome would narrow or reject
the proposed MTT string encoding rather than being rephrased as success.

\section{Questions the slogan often hides}

\subsection{Is the worldsheet literally spacetime?}

No.  It is the two-dimensional surface traced by the string.  Its coordinates
and RG scale are not the coordinates and time of target spacetime.

\subsection{Does vanishing of one beta function suffice?}

No.  The full background, anomaly, ghost, modular, GSO, factorization, and
global data matter.  In heterotic theory the bundle and Green--Schwarz sector
are indispensable.

\subsection{Does this quantize gravity in the usual sense?}

Perturbative closed strings contain a massless spin-two excitation and provide
quantum gravitational amplitudes in their regime.  The MTT interpretation
explored here is different: it asks whether quantized upper dynamics and
projection produce the gravitational effective medium and its stress response.
That proposal still needs the selected upper action and quantum transport.

\subsection{Are GR and string theory equally fundamental?}

The beta-function bridge does not decide ontology.  It establishes a
controlled relation between descriptions.  MTT's preprojection proposal would
place both downstream of an upper source, but only after that source is
constructed.

\section{Conclusion}

General Relativity appears in perturbative string theory through a beautiful
but qualified chain.  Target fields are worldsheet couplings; quantum Weyl
consistency constrains those couplings; the constraints agree with target
effective equations; and a restricted low-energy compactified corner yields
Einstein gravity.

MTT may eventually explain why the worldsheet and spacetime diagnostics agree
by deriving both from one selected preprojection structure.  The technical
conditions for such an explanation are now stated in the companion
diagnostic-square paper.  The physical q79 source is not complete: five
worldsheet rows are available, two are partial, and five are open.

The right conclusion is therefore neither ``string theory and GR are the same''
nor ``the relation is only an analogy.''  The standard perturbative bridge is
real.  The MTT common-source explanation is a precise, testable, unfinished
program.

\begin{thebibliography}{99}

\bibitem{Friedan1980}
D.~Friedan,
``Nonlinear models in \(2+\epsilon\) dimensions,''
\emph{Physical Review Letters} \textbf{45} (1980), 1057--1060.
\url{https://doi.org/10.1103/PhysRevLett.45.1057}

\bibitem{CallanEtAl1985}
C.~G. Callan, D.~Friedan, E.~J. Martinec, and M.~J. Perry,
``Strings in background fields,''
\emph{Nuclear Physics B} \textbf{262} (1985), 593--609.
\url{https://doi.org/10.1016/0550-3213(85)90506-1}

\bibitem{MetsaevTseytlin1987}
R.~R. Metsaev and A.~A. Tseytlin,
``Order \(\alpha'\) (two-loop) equivalence of the string equations of motion
and the sigma-model Weyl invariance conditions,''
\emph{Nuclear Physics B} \textbf{293} (1987), 385--419.
\url{https://doi.org/10.1016/0550-3213(87)90077-0}

\bibitem{HullTownsend1988}
C.~M. Hull and P.~K. Townsend,
``String effective actions from sigma-model conformal anomalies,''
\emph{Nuclear Physics B} \textbf{301} (1988), 197--223.
\url{https://doi.org/10.1016/0550-3213(88)90692-5}

\bibitem{NeroSquare2026}
P.~Nero,
\emph{Worldsheet and Spacetime Consistency as a Conditional Diagnostic
Square},
Version 2, 2026.
\url{https://doi.org/10.5281/zenodo.21719524}

\bibitem{NeroProjectionString}
P.~Nero,
\emph{A Projection-First Reframing of String Theory: Conditional Encodings,
Worldsheet Gates, and the q79 Boundary},
Version 2, 2026.

\end{thebibliography}
\end{document}
